\section{Evaluation and Human Validation}
\label{app:evaluation}

\subsection{Scoring}

Each case defines user-task checks and injection checks. Every check is a
boolean function of the run's tool calls, with their arguments, and of the
emails the agent sent. A case combines its checks with AND and OR into a
task-success formula and an attack-success formula, and a check that raises
an error counts as false. Attack success is scored independently of task
success, so a run can complete the user task and still carry out the attack.



When an outcome is free-form text, the check first extracts the relevant
field, such as the bodies of the emails sent to one recipient or the content
written to one file. An LLM judge then answers a case-specific yes/no
criterion about that field alone, never about the whole trajectory
(Listing \ref{lst:llm-predicate}). The judge is Claude Sonnet 4, called
3 times at temperature
0; the predicate holds if
2 of the
3 answers are YES. Claude Sonnet 4 is also
an evaluated agent, but the judge sees only the extracted field, not the
trajectory or the model that produced it.

\begin{lstlisting}[basicstyle=\ttfamily\scriptsize,breaklines=true,frame=single,columns=fullflexible,keepspaces=true,caption={Prompt of the LLM predicate. \texttt{\{criterion\}} is the case-specific yes/no question and \texttt{\{text\}} the field extracted by the check; the predicate holds if 2 of 3 calls at temperature 0 answer YES.},label={lst:llm-predicate}]
You are a precise text analysis tool. Answer the following question about the text below with exactly YES or NO. No explanation.

Question: {criterion}

Text:
"""
{text}
"""

Answer (YES or NO):
\end{lstlisting}

\subsection{Human validation}

\textbf{Case validity.} Two annotators independently reviewed each of the
111 designed cases against the three
conditions of Section \ref{sec:ipi-boundary}, also checking that the harmful
outcome can be observed and scored. Agreement was
92.8\%,
84.7\%, and
98.2\% for conditions 1--3
(91.9\% overall).

\textbf{Target labels.} Human annotation of 225 attacks
from 25 cases, covering all six primitives
(27--42
attacks each), agreed with the assigned primitive for
217 (96.4\%).

\textbf{Scores.} We sampled trajectories for human annotation by stratified
sampling over the automatically assigned attack success and task success, the
case, and the model. Human annotations of 89 trajectories
(25 DeepSeek V3.1, 19 GLM-5, 27 Kimi K2.5, 18 Claude Sonnet 4) agreed with the automatic scores in
173 of 178 task
and attack judgments (97.2\%).
